% Source: https://artofproblemsolving.com/community/c3h1514698_line_geometry
% Stable identity: aops-topic-1514698-q01
Let $K=(k,0)$. Since $AK=3BK>BK$, the point $K$ is closer to $B$ than to $A$, so it cannot lie to the left of $A$ (there $AK<BK$). Being outside the segment, it lies to the right of $B$: $k>2$, so $AK=k+3$ and $BK=k-2$. Then
\[
k+3=3(k-2)\qquad\Longrightarrow\qquad k=\dfrac{9}{2} .
\]

\begin{asy}
usepackage("amsmath");
size(6.5cm);
defaultpen(fontsize(10pt));
draw((-4.5,0)--(6,0),Arrow(TeXHead));
for(int i=-4;i<=5;++i) draw((i,-0.08)--(i,0.08));
pair A=(-3,0),B=(2,0),K=(4.5,0);
draw(A--B,linewidth(1.6));
dot(A);dot(B);dot(K,red);
label("$A(-3,0)$",A,N);label("$B(2,0)$",B,N);label("$K\left(\dfrac{9}{2},0\right)$",K,N,red);
draw((-3,-0.5)--(4.5,-0.5),Arrows(TeXHead));label("$AK=\dfrac{15}{2}$",(-0.6,-0.5),S);
draw((2,-1.6)--(4.5,-1.6),Arrows(TeXHead));label("$BK=\dfrac{5}{2}$",(3.25,-1.6),S);
label("\mbox{}",(min(currentpicture,user=true)+max(currentpicture,user=true))/2,black); // keeps the text after the figure black when LaTeX typesets the labels inline
\end{asy}

Check: $AK=\dfrac{15}{2}$ and $BK=\dfrac{5}{2}$. So $K=\boxed{\left(\dfrac{9}{2},0\right)}$.

(The ratio of lengths is positive. Ratios can be given a sign when \emph{directed} segments are used: then $\dfrac{\overline{AK}}{\overline{KB}}$ is negative for points outside the segment. The point inside the segment with $AK=3BK$ is $\left(\dfrac{3}{4},0\right)$, excluded here.)
